% 135-02(135쪽) — 함수 $y=\sin x\,(x\ge 0)$ 의 그래프와 직선 $y=k\,(0<k<1)$ 의 교점 $x_1\sim x_5$.
% 스캔 화질이 낮아 TikZ 로 다시 그림. 원문에 있는 라벨(O · x_1~x_5 · y=sin x · y=k)만 두고
% 눈금값(π, 2π …)은 원문처럼 넣지 않는다 — 교점의 값이 그림에서 읽히지 않게.
\begin{tikzpicture}[line width=0.6pt, x=0.44cm, y=0.96cm, >=stealth]
  \draw[->, line width=0.7pt] (-1.0,0) -- (15.5,0) node[below=1pt, font=\footnotesize] {$x$};
  \draw[->, line width=0.7pt] (0,-1.45) -- (0,1.5) node[left=1pt, font=\footnotesize] {$y$};
  \draw[line width=0.5pt] (-0.7,0.55) -- (14.55,0.55) node[right=1pt, font=\footnotesize] {$y=k$};
  \draw[domain=0:13.95, samples=220, smooth] plot (\x, {sin(\x r)});
  \foreach \xv/\lab in {0.582/x_1, 2.559/x_2, 6.866/x_3, 8.842/x_4, 13.149/x_5} {
    \draw[densely dotted, line width=0.35pt] (\xv,0.04) -- (\xv,0.54);
    \node[below=-1.5pt, font=\footnotesize] at (\xv,0) {$\lab$};
  }
  \node[below left=-2.5pt and -2.5pt, font=\footnotesize] at (0,0) {$\pt{O}$};
  \draw[->, line width=0.35pt] (8.5,1.2) -- (7.75,0.92);
  \node[anchor=west, font=\footnotesize] at (8.4,1.3) {$y=\sin x$};
\end{tikzpicture}
